\documentclass[11pt,a4paper]{article}
\usepackage[margin=22mm]{geometry}
\usepackage{kotex}
\usepackage{amsmath,amssymb,mathtools,bm}
\usepackage{booktabs,longtable,array,tabularx}
\usepackage{enumitem}
\usepackage{hyperref}
\usepackage{amsthm}
\usepackage{mdframed}
\usepackage{fancyhdr}
\usepackage{lastpage}
\usepackage{setspace}
\setstretch{1.13}
\setlength{\parskip}{0.4em}
\setlength{\parindent}{0pt}
\setlength{\headheight}{14pt}
\hypersetup{colorlinks=true, linkcolor=blue, urlcolor=blue, citecolor=blue,
  pdftitle={Graphite ICA tail: barrier-spectrum kernel-integral theory}}
\pagestyle{fancy}\fancyhf{}
\lhead{Graphite ICA tail — Chapter 1}\rhead{\thepage/\pageref{LastPage}}
\renewcommand{\headrulewidth}{0.4pt}

\newtheorem*{groundingbox}{Grounding}
\newtheorem*{boundbox}{Validity bound}
\newtheorem*{flagbox}{Flagged (미확정)}
\newtheorem*{uservb}{사용자 verbatim}

\newcommand{\dd}{\mathrm{d}}
\newcommand{\eff}{\mathrm{eff}}
\newcommand{\eq}{\mathrm{eq}}
\newcommand{\app}{\mathrm{app}}
\newcommand{\dimedrive}{\mathrm{drive}}
\newcommand{\tot}{\mathrm{tot}}
\newcommand{\bg}{\mathrm{bg}}
\newcommand{\cell}{\mathrm{cell}}
\newcommand{\ext}{\mathrm{ext}}
\newcommand{\simplecase}{\mathrm{simple}}
\newcommand{\AL}[1]{\textsf{[AL-#1]}}

\title{\textbf{리튬이온전지 graphite anode 의 ICA(\,$dQ/dV$\,) tail 거동:\\
effective barrier 와 relaxation-length spectrum kernel-integral 이론}\\[0.4em]
\large Chapter 1}
\author{Project\_Anode\_Fit}
\date{}

\begin{document}
\maketitle

\begin{center}\begin{minipage}{0.93\textwidth}
\small \textbf{Abstract.} 리튬이온전지 graphite anode 의 incremental capacity analysis
(ICA), $dQ/dV$ 곡선에서 staging transition peak 뒤의 \emph{post-peak tail} 은 temperature
가 낮을수록 길게 늘어지고(stretched) 높을수록 짧게 끝난다. 본 Chapter 1 은 이 tail 을
\textbf{단일 relaxation length 가 아니라 relaxation-length spectrum 의 kernel integral}
로 해석한다. 출발점은 \textbf{activation free-energy barrier 가 state 를 $0\!\to\!1$ 로
점프시키는 threshold 가 아니라 rate constant 를 정하는 hidden variable} 이라는 점이다.
각 local mode 는 effective barrier $\Delta G_\eff = \Delta G_a-\chi\mathcal A$ 가 정하는
Eyring-type rate 로 equilibrium target 을 향해 smooth relaxation 하며, 하나의 exponential
kernel ($L_q=v_q/k=|I|/(Q_\cell k)$) 을 만든다. 관측 tail 은 barrier distribution $\rho_G$ 가 \emph{지수적
barrier$\to$length 매핑}을 통해 만드는 relaxation-length spectrum $A_L$ 에 대한 kernel
integral 이다. Self-consistent integral closure 는 물리 core 와 분리한다. \textbf{Plan A}는
사용자 박사 연구의 Refs~\cite{lee2011,son2013}에서 온 Fredholm / propagator / ratio-substitution
method가 현재 graphite ICA 구조에 정합할 때만 쓰는 candidate analytic closure 이다.
\textbf{Plan B}는 계산 절차가 아니라, local residual ODE 와 relaxation-length spectrum
kernel integral로 구성되는 conservative theoretical formulation 이다. Plan A가 검증되지 않아도
Plan B만으로 Chapter 1의 물리 논리, 즉 low temperature 에서 large-$L_q$ weight가 커지고
potential assistance가 short-$L_q$ shift를 만든다는 설명은 유지된다. Equilibrium target 의 함수
형태는 특정하지 않고, grounded lattice-gas를 특수해로 둔다. $V_n$ 은 외부 lookup 이 아니라
charge conservation 으로 implicit 결정된다. 모든 assumption 은 Assumption Ledger (AL-\#)로
위치를 표시하고, tail 을 effective barrier에 귀속하기 전에 polarization, transport,
heterogeneous width, charge-transfer co-linearity를 분리하는 falsification protocol을 둔다.
\end{minipage}\end{center}

\tableofcontents
\newpage

% ====================================================================
\section{Motivation 과 deliverable}\label{sec:intro}
% ====================================================================

\subsection{사용자 의도 (verbatim)}\label{sec:verbatim}
\begin{uservb}
``LIB의 ICA 분석에서 온도와 현시점의 전위 상태에 따라서 그 그래프의 개형이 특히
피크부분이 상변이에 의해서 나타나는데 그 상변이의 의해 나타나는 피크 꼬리의 개형이
늘어지느냐 빨리 떨어지느냐가 달라지는것을 확인했다. 특히 온도가 낮을때는 꼬리가
길고, 온도가 높으면 그 꼬리가 상대적으로 짧게 끝나는 현상이 관측되었어. 원래라면
그렇게 온도에 대해 약간 가우시안 피크 형상이 나타나고 더이상 진행이 안될 것으로
생각되는데 그 온도에 의한 배리어 외에 추가적인 극판 자체 전위에 따른 배리어를
낮추는 효과가 있다고 봤고 그 유효 배리어의 논리를 확정하고 그 것을 이용해 피크의
거동을 해석하는 피팅하는데 쓸수 있는 논리 식을 구하는게 이 작업의 핵심이야.''
\end{uservb}

\textbf{기준 해석.} (1) ``가우시안''은 \emph{sharp peak} 를 가리키는
예시이며 equilibrium 형태를 Gaussian 으로 확정한 것이 아니다; 형태는 Assumption Ledger 와 data
로 결정한다. (2) ``activation energy barrier''는 본 이론의 핵심으로 \emph{유지}하되, ``특정
지점을 넘으면 $0\!\to\!1$ 로 갑자기 점프''하는 step-function 식 \emph{비약은 금지} (모든
transition 은 smooth). 본 chapter 는 이 두 기준을 절대 기준으로 한다.

\subsection{관측 사실 (observation)}\label{sec:observation}
\begin{enumerate}[label=\textbf{(O\arabic*)}]
\item \textbf{Peak 출현.} graphite staging transition 영역에서 $dQ/dV_{n,\app}$ 가 sharp
  peak 를 갖는다 (staging sequence 가 $N_p\approx3\sim4$ effective transition 으로 fusion).
\item \textbf{Peak tail 의 stretched temperature 의존.} 각 peak 의 tail 이 $T$ 가
  낮을수록 길게 늘어지고(흔히 단일 exponential 보다 \emph{stretched}) 높을수록 짧게
  끝난다. \emph{단일 equilibrium peak width 만으로는} 이 $T$/potential 의존 tail 을
  설명하기 어렵다.
\end{enumerate}
본 Chapter 의 1차 초점은 (O2). peak area 정합은 fitting 단계로 이관 (P1).

\subsection{가설 (hypothesis)}\label{sec:hypothesis}
\textbf{(H1)} equilibrium target $\theta_{\eq,j}(V_n,T)$ 는 $V_n$ 에 대해 smooth 하게
변한다 ($0\!\to\!1$ 급점프 아님; \S\ref{sec:equilibrium}).

\textbf{(H2)} 각 local mode 의 진행 속도는 activation barrier $\Delta G_{a}$ 를 surmount 하는
thermally activated process 이다 (\AL{1}).

\textbf{(H3, ★)} electrode potential $V_{n,\dimedrive}$ 가 그 barrier 를 $\chi\,\mathcal A$
만큼 낮춘다: $\Delta G_\eff = \Delta G_a - \chi\mathcal A$, $\mathcal A=s_\phi F(V_{n,\dimedrive}-U)$
(\AL{2}; 정식 \S\ref{sec:barrier}).

\textbf{(H4)} barrier 는 rate constant $k=\nu\exp(-\Delta G_\eff/RT)$ 를 정하고 (\AL{1}),
$\theta_j$ 는 target $\theta_{\eq,j}$ 를 finite rate 로 따라가 lag 를 만든다 (\AL{5}).
하나의 mode 는 normalized charge coordinate 에서 length $L_q=v_q/k=|I|/(Q_\cell k)$ 의 \emph{단일 exponential kernel} 을 만든다.

\textbf{(H5, ★ 핵심)} 실제 electrode 는 단일 barrier 가 아니라 \emph{barrier
distribution} $\rho_G$ 를 가진다. barrier$\to$length 의 \emph{지수} 매핑이 그 분산을
크게 확대하여 \emph{relaxation-length spectrum} $A_L$ 을 만들고, 관측 tail 은 그 spectrum
의 \textbf{kernel integral} 이다 (\S\ref{sec:spectrum}--\ref{sec:kernel}). 따라서 관측
tail 은 일반적으로 stretched 이다.

\textbf{(H6)} $T$ 가 낮으면 spectrum 이 large-$L_q$ 로 이동해 tail 이 길어지고, $T$ 가
높거나 potential assistance 가 크면 short-$L_q$ 로 이동해 tail 이 짧아진다 (O2; \S\ref{sec:tailT}).

\subsection{Deliverable}\label{sec:deliverable}
(H1)$\sim$(H6) 을 grounded derivation 으로 확정하고, ICA tail 의 $T$/potential 의존을
설명하는 \emph{논리식}(spectrum kernel integral 과 conservative theoretical formulation)을 도출한다. 단 (i) 모든 assumption 은 AL-\# grounding, (ii) FLAGGED 는 established
로 쓰지 않음, (iii) tail 의 barrier 귀속 전에 competing source 의 falsification protocol
제시 + barrier distribution 의 non-unique 역산 회피, (iv) quantitative fitting implementation 은 Chapter 1 범위 외 (P1).

\subsection{사용자 박사 연구의 위치 (Refs 6/7) = candidate analytic closure}\label{sec:phd}
spectrum kernel integral 과 self-consistent relaxation 은 미지 함수가 적분 안팎에 동시에
나타나는 integral equation 을 만든다. 사용자 박사 연구 \cite{lee2011,son2013} 의
\emph{propagator / ratio-substitution} 기법은 이런 적분방정식에 대해 Plan A candidate analytic
closure 를 제공할 수 있는 수학적 도구다. 그러나 본 Chapter 의 load-bearing 물리 core 는
activation barrier $\to$ rate constant $\to$ relaxation length $\to$ spectrum kernel integral
chain 이며, Refs 6/7 은 이 core 를 편리하게 닫는 후보 방법으로만 둔다. \textbf{Validation gate
(\AL{7})}: 그 기법은 \emph{Fredholm} integral equation of the second kind 를 대상으로
하므로, 본 형식이 \emph{Volterra} 인지 Fredholm 인지에 따른 적용성을 검증 대상(VG-2)
으로 둔다.

\subsection{Grounded spine 요약 (S1--S14, synthesis)}\label{sec:spine}
\begin{longtable}{p{0.06\textwidth} p{0.66\textwidth} p{0.2\textwidth}}
\toprule ID & 골격 & 상태 [AL] \\ \midrule \endhead
S1 & charge conservation $\Rightarrow$ implicit $V_n$ & GROUNDED \AL{9} \\
S2 & equilibrium target $\theta_\eq$ (일반형, smooth; lattice-gas 특수해) & GROUNDED \AL{3a} \\
S3 & broadening 형태 미확정 (homog.\ interaction / heterog.\ disorder); S2/S8 에 흡수 & FLAGGED \AL{3b} \\
S4 & activation free energy $\Delta G_a(T)$ & GROUNDED \AL{1} \\
S5 & ★ effective barrier $\Delta G_\eff=\Delta G_a-\chi\mathcal A$ & BOUNDED \AL{2} \\
S6 & single-mode relaxation $\dot\theta=k(\theta_\eq-\theta)$ & BOUNDED \AL{5} \\
S7 & single-mode exponential kernel $L_q=v_q/k=|I|/(Q_\cell k)$; $L_Q=Q_\cell L_q$ & GROUNDED \AL{1} \\
S8 & ★ barrier dist.\ $\rho_G\to$ relaxation-length spectrum $A_L$ & BOUNDED \AL{6} \\
S9 & ★ kernel integral for observed tail & BOUNDED \AL{6} \\
S10 & ★ closure: Plan B (conservative theoretical formulation, 항상 보존) + Plan A (Refs 6/7 candidate analytic closure, gated) + R1-R5 governance & BOUNDED+GATED \AL{7} \\
S11 & ICA mapping $dQ_\ext/dV_n$ & GROUNDED \AL{9} \\
S12 & ★ spectrum 의 $T/\psi$ shift = stretched tail & FLAGGED novel \AL{8} \\
S13 & ★ falsification + non-uniqueness 처리 & method \AL{10} \\
S14 & later fitting interface + identifiability + validation & GROUNDED \AL{10} \\
\bottomrule
\end{longtable}
\textbf{Validation rule.} 모든 assumption 은 AL-\# 인용. FLAGGED 는 established 로 사용 금지; BOUNDED
는 유효범위 병기. step-function·$0\!\to\!1$ 급점프·$\max$·$\min$·Heaviside·hard-switch
정의식 금지 (smooth only).
\subsection{Compact Assumption Ledger}\label{sec:assumption_ledger}
\begin{longtable}{p{0.11\textwidth} p{0.17\textwidth} p{0.60\textwidth}}
\toprule ID & 상태 & 본 Chapter 에서의 의미 \\ \midrule \endhead
AL-1 & GROUNDED & transition-state / Eyring 관점: activation free-energy barrier 는 rate constant 를 정한다. \\
AL-2 & BOUNDED & BEP/Marcus small-driving-force 영역에서 electrode potential assistance 는 effective barrier 를 smooth 하게 낮춘다. \\
AL-3a & GROUNDED & equilibrium target 의 grounded 특수해는 lattice-gas / regular-solution chemical potential 에서 온다. \\
AL-3b & FLAGGED & erf/Gaussian-cumulative target 은 graphite thermodynamic derivation 이 확정되기 전까지 empirical ansatz 다. \\
AL-4 & GROUNDED & graphite staging transition 은 ICA peak 의 thermodynamic origin 이다. \\
AL-5 & BOUNDED & local mode 는 near-equilibrium first-order relaxation 으로 선형화할 수 있다. \\
AL-6 & BOUNDED & barrier distribution 에서 relaxation-length spectrum 으로 가는 역매핑은 non-unique 이므로 forward prediction 만 허용한다. \\
AL-7 & GATED & Refs 6/7 Fredholm / propagator / ratio-substitution method 는 Plan A candidate closure 이며, 적용성 검증 전에는 core assumption 이 아니다. \\
AL-8 & FLAGGED & low-$T$ long tail 과 potential-assisted short-$L_q$ shift 를 kinetic-spectrum signature 로 해석하는 부분은 novel hypothesis 이다. \\
AL-9 & GROUNDED & internal anode potential $V_n$ 은 charge conservation 의 implicit root 이며 외부 lookup 이 아니다. \\
AL-10 & METHOD & competing mechanisms, identifiability, and falsification protocol 을 만족해야 barrier-spectrum 귀속을 주장할 수 있다. \\
\bottomrule
\end{longtable}

% ====================================================================
\section{Notation 과 convention}\label{sec:notation}
% ====================================================================
\begin{longtable}{p{0.16\textwidth} p{0.14\textwidth} p{0.62\textwidth}}
\toprule 기호 & 단위 & 의미 \\ \midrule \endhead
$q$ & --- & 정규화 진행 좌표 $q=Q_\ext/Q_\cell$; 본 Chapter 의 primary coordinate \\
$Q_\ext$ & C & dimensional external charge; $v_Q\equiv dQ_\ext/dt=|I|>0$ \\
$Q_\cell$ & C & 기준 cell capacity; $q=Q_\ext/Q_\cell$ 의 scale \\
$v_q$ & 1/s & normalized scan speed $v_q\equiv dq/dt=v_Q/Q_\cell=|I|/Q_\cell$ \\
$T,F,R,k_B,h$ & --- & temperature, Faraday, gas const, Boltzmann, Planck \\
$V_n$ & V & internal anode (analysis) potential; charge conservation 의 root \\
$V_{n,\app}$ & V & apparent(measured) potential $=V_n+s_I|I|R_n$ \\
$V_{n,\dimedrive}$ & V & driving potential (rate/barrier 인수); 기본 $=V_n$, reduced $\approx V_{n,\app}$ \\
$R_n(q,T,|I|)$ & $\Omega$ & 유효 polarization 계수 \\
$j,\,N_p$ & --- & effective transition index, 개수 \\
$\theta_j,\ \theta_{\eq,j}$ & --- & local mode 진행률, equilibrium target ($0\!\le\!\theta\!\le\!1$) \\
$r_j$ & --- & kinetic lag $=\theta_{\eq,j}-\theta_j$ \\
$U_j(T),w_j(T)$ & V & equilibrium 중심 전위, isotherm width \\
$\Delta H_{a},\Delta S_{a},\Delta G_a$ & J/mol & activation enthalpy/entropy/free energy \\
$\mathcal A_j$ & J/mol & driving force $=s_{\phi,j}F(V_{n,\dimedrive}-U_j)$ \\
$\chi_j$ & --- & transfer coefficient (BEP), $0\le\chi_j\le1$ \\
$W_\psi$ & J/mol & potential barrier-lowering work $=\chi_j\mathcal A_j$ \\
$\Delta G_{\eff}$ & J/mol & ★ effective barrier $=\Delta G_a-\chi\mathcal A=G-W_\psi$ \\
$G$ & J/mol & local activation barrier (mode 별; hidden rate variable) \\
$\nu(T),\kappa(T),k_0(T)$ & 1/s,--,1/s & Eyring prefactor $\nu=(k_BT/h)\kappa$ \\
$k$ & 1/s & rate constant \\
$L_q$ & --- & normalized local relaxation length in $q$ coordinate, $L_q=v_q/k=|I|/(Q_\cell k)$ \\
$L_Q$ & C & dimensional charge-axis relaxation length, $L_Q=v_Q/k=Q_\cell L_q$; dimensional $Q_\ext$ 식에서만 사용 \\
$\rho_G(G;T,\psi)$ & 1/(J/mol) & activation barrier distribution \\
$A_L(\lambda;T,\psi)$ & 1 & relaxation-length spectral \emph{density} over normalized length variable $\lambda=L_q$; $\int A_L(\lambda)\,d\lambda$ 가 weight \\
$\Theta$ & --- & 전체 effective phase progress; $Q_p$ = 그 capacity scale \\
$C_\bg(V_n,T)$ & C/V & background chemical capacitance \\
\bottomrule
\end{longtable}
\begin{groundingbox}
\textbf{Smooth-only.} 모든 정의식은 smooth. 금지 모델의 대표 예: $\theta=H(G_c-G)$
(Heaviside threshold completion) — barrier 를 넘으면 state 가 급점프한다는 모델. 본 chapter
는 이를 채택하지 않는다 (smooth-only convention). 단위: $\Delta G$ J/mol, $\Delta G/RT$
무차원, $k,\nu$ 1/s, $V$ V.
\end{groundingbox}

% ====================================================================
\section{Graphite staging 과 effective transition}\label{sec:staging}
% ====================================================================
Graphite lithiation 은 stage 4 $\to$ 3 $\to$ 2L $\to$ 2 $\to$ 1 sequence 로 진행하며 각
transition 이 $dQ_\ext/dV_n$ 에 peak 를 남긴다 \cite{dahn1991,ohzuku1993,funabiki1999} (\AL{4}).
인접 stage 의 중심 전위 차가 width 이하이면 fusion 되어 observable peak 수가 $N_p\approx
3\sim4$ 가 된다. 본 Chapter 는 이를 effective transition $j=1,\dots,N_p$ 로 둔다. 단 각
effective transition 내부에도 particle/domain/boundary 다양성에 의한 \emph{local mode 분포}
가 존재하며, 이것이 \S\ref{sec:spectrum} 의 spectrum 의 물리적 근원이다.

% ====================================================================
\section{Charge conservation 과 implicit $V_n$}\label{sec:charge}
% ====================================================================
\textbf{물리.} anode 내부 charge = 외부 charge (conservation). 이것이 $V_n$ 을 결정한다
(외부 함수 아님; \AL{9}).

\textbf{Derivation D1.} background 와 phase progress 기여의 합
\begin{equation}
Q_\cell\,q = Q_\bg(V_n,T) + Q_p\,\Theta(V_n,q,T),\qquad \Theta=\frac1{N_p}\sum_j w^\Theta_j\theta_j,
\label{eq:charge_balance}
\end{equation}
를 $V_n$ 에 대해 풀어 $V_n$ 을 얻는다 (single-particle/porous-electrode 의 implicit $\mu$
관계와 동형 \cite{doyle1993,newman}). unique smooth 해 조건:
\begin{equation}
\frac{\partial Q_\bg}{\partial V_n}\ge\epsilon_Q>0.
\label{eq:monotone}
\end{equation}
$\theta_j=\theta_{\eq,j}$ 이면 그 해를 $V_n=V_{n,\mathrm{OCV}}(q,T)$ 로 정의 — 외부 OCV
lookup 은 이 charge-balance 의 equilibrium 특수해다.
\textbf{Derivation D1b (differential charge balance).} 식~\eqref{eq:charge_balance} 를
solution path 를 따라 $q$ 로 미분하면
\begin{equation}
Q_\cell
=
C_\bg(V_n,T)\frac{\dd V_n}{\dd q}
+Q_p\frac{\dd\Theta}{\dd q},
\qquad
C_\bg\equiv \frac{\partial Q_\bg}{\partial V_n}.
\label{eq:charge_balance_differential}
\end{equation}
따라서
\begin{equation}
\frac{\dd V_n}{\dd q}
=
\frac{Q_\cell-Q_p\,\dd\Theta/\dd q}{C_\bg(V_n,T)}.
\label{eq:dvdq_charge}
\end{equation}
이 식은 later ICA mapping 의 출발점이며, $\dd\Theta/\dd q$ 에 tail contribution 이 들어간다.

% ====================================================================
\section{Equilibrium target $\theta_\eq$ (일반형, smooth)}\label{sec:equilibrium}
% ====================================================================
\textbf{보수적 입장.} 본 Chapter 는 equilibrium target 의 함수 형태를 특정하지 않는다.
요구 조건은 (i) smooth, (ii) $0\le\theta_\eq\le1$, (iii) $V_n$ 의 단조 함수뿐이다:
\begin{equation}
\theta_{\eq,j}=\theta_{\eq,j}(V_n,T),\qquad 0\le\theta_{\eq,j}\le1.
\label{eq:thetae}
\end{equation}
(사용자 정정: ``가우시안''은 sharp peak 예시일 뿐 형태 확정 아님 → 특정 형태 강제 X.)

\textbf{Grounded 특수해 (lattice-gas, \AL{3a}).} 명시적 형태가 필요하면, graphite
intercalation 의 grounded thermodynamics 인 lattice-gas / regular-solution chemical
potential \cite{mckinnon1983,bazant2013} 을 쓴다:
\begin{equation}
\mu_j(\theta_j,T)
=
\mu_j^0(T)
+RT\ln\frac{\theta_j}{1-\theta_j}
+\Omega_j(1-2\theta_j).
\label{eq:mu_regular_solution}
\end{equation}
equilibrium 에서는 electrochemical potential balance 가 성립하므로
\begin{equation}
\mu_j(\theta_{\eq,j},T)=s_{\phi,j}F\big(V_n-U_j^0(T)\big).
\label{eq:mu_equilibrium_condition}
\end{equation}
ideal limit $\Omega_j=0$ 에서 식~\eqref{eq:mu_equilibrium_condition} 을 $\theta_{\eq,j}$ 에 대해
풀면
\begin{equation}
\theta_{\eq,j}
=
\big[1+\exp(-s_{\phi,j}(V_n-U_j)/w_j)\big]^{-1},
\qquad
w_j=\frac{RT}{F},
\label{eq:latticegas}
\end{equation}
이고, 따라서
\begin{equation}
\frac{\partial\theta_{\eq,j}}{\partial V_n}
=
\frac{s_{\phi,j}}{w_j}\theta_{\eq,j}(1-\theta_{\eq,j}),
\qquad
\left|\frac{\partial\theta_{\eq,j}}{\partial V_n}\right|_\mathrm{max}
=
\frac{1}{4w_j}.
\label{eq:thetaeq_derivative}
\end{equation}
즉 homogeneous lattice-gas target 의 voltage width 는 $w_j=RT/F$ 로 temperature 에 비례한다.
저온에서는 equilibrium peak 가 더 좁아지는 방향이므로, 관측된 \emph{low-$T$ long tail} 을
homogeneous equilibrium width 만으로 설명할 수 없다. $\Omega_j\ne0$ 이면 width/형태가 보정되지만
target 은 여전히 smooth 이며 $0\!\to\!1$ 급점프가 아니다.
\begin{flagbox}
\textbf{형태는 data 결정 (VG-1).} erf/Gaussian-cumulative isotherm 은 graphite
intercalation 에 대한 thermodynamic derivation 이 문헌에 없으므로 (organic-semiconductor
Gaussian Disorder Model \cite{baessler1993} 과의 analogy + 경험 peak-fit 수준) FLAGGED
ansatz 로만 둔다 (\AL{3b}). 실제 형태는 저속 OCV peak 의 near-peak lineshape 으로 판별
(\S\ref{sec:falsify}). 본 Chapter 의 tail 결론은 $\theta_\eq$ 의 \emph{구체 형태에
의존하지 않는다} — post-peak 에서 $d\theta_\eq/dq\simeq0$ 만 쓰기 때문 (\S\ref{sec:kernel}).
\end{flagbox}

% ====================================================================
\section{★ Effective barrier $\Delta G_\eff$}\label{sec:barrier}
% ====================================================================
\textbf{Derivation D2 (activation free energy, \AL{1}).} transition-state theory
\cite{eyring1935,eyringchemrev1935} 에서 $\Delta G_a(T)=\Delta H_a-T\Delta S_a$.

\textbf{Derivation D3 (potential lowering, BEP; \AL{2}).} 외부 driving force 가 barrier 를
부분적으로 낮춘다는 것은 reaction-coordinate 상 transition state 위치(Hammond)에 따른
비율 $\chi$ 로 표현된다 (Brønsted--Evans--Polanyi \cite{evanspolanyi1938}; Butler--Volmer
transfer coefficient \cite{bardfaulkner}). driving force $\mathcal A_j=s_{\phi,j}F(V_{n,\dimedrive}
-U_j)$ 에 대해 lowering work $W_\psi=\chi_j\mathcal A_j$, 따라서
\begin{equation}
\boxed{\;\Delta G_{\eff}=\Delta G_a-\chi_j\,\mathcal A_j\;=\;G-W_\psi.\;}
\label{eq:Geff}
\end{equation}
여기서 $G\equiv\Delta G_a$ (mode 의 intrinsic barrier; \S\ref{sec:spectrum} 에서 분포),
$W_\psi=\chi_j\mathcal A_j$ (uniform potential lowering). $\mathcal A_j>0$ 이면 barrier 가
낮아진다 — 사용자 verbatim 의 ``극판 전위에 따른 배리어를 낮추는 효과''.
\begin{boundbox}
\textbf{Marcus bound (\AL{2}).} 선형 lowering 은 small driving force 의 leading-order
근사다. Marcus 이론 \cite{marcus1956} 은 barrier 가 driving force 의 2차 함수임을 보이며
큰 $|\mathcal A_j|$ 에서 곡률/inverted region 으로 선형 근사가 깨진다. 본 선형식은
$V_{n,\dimedrive}\approx U_j$ 인 tail 영역에서 유효; 전역 정당화 아님.
\end{boundbox}
\textbf{$\Delta G_\eff$ 가 작아지는 영역 (step 금지).} $W_\psi$ 가 $G$ 에 근접해 activation
이 더 이상 rate-limiting 이 아니면, 식~\eqref{eq:Geff} 를 $\max(\cdot,0)$/clip 으로 고치지
\emph{않는다}. 그 영역에서는 attempt frequency, site availability, phase-boundary motion,
diffusion, transport, current conservation 같은 \emph{다른 physical bottleneck} 이 rate 를
제한한다 (\S\ref{sec:falsify}).

% ====================================================================
\section{Rate constant 과 single-mode exponential kernel}\label{sec:rate}
% ====================================================================
\textbf{Derivation D4 (Eyring rate, \AL{1}).}
\begin{equation}
k(G,T,\psi)=k_0(T)\exp\!\Big[-\frac{G-W_\psi}{RT}\Big],\qquad k_0(T)=\frac{k_BT}{h}\kappa(T).
\label{eq:rate}
\end{equation}
prefactor 는 $k_BT/h$ (Eyring) 이며 임의 상수가 아니다.

\textbf{Derivation D5 (single-mode relaxation, \AL{5}).} 한 local mode 의 lag 를
$r=\theta_\eq-\theta$ 로 둔다. forward/backward elementary rate constant 를 $k_+$, $k_-$ 라 쓰면
\begin{equation}
\dot\theta
=
k_+(1-\theta)-k_-\theta.
\label{eq:forward_backward_flux}
\end{equation}
stationary condition $\dot\theta=0$ 에서
\begin{equation}
\theta_\eq=\frac{k_+}{k_++k_-}.
\label{eq:thetaeq_rates}
\end{equation}
식~\eqref{eq:forward_backward_flux} 에 식~\eqref{eq:thetaeq_rates} 를 대입하면
\begin{equation}
\dot\theta
=
(k_++k_-)(\theta_\eq-\theta)
\equiv
k(\theta_\eq-\theta).
\label{eq:first_order_relaxation}
\end{equation}
즉 $\theta_\eq$ 는 thermodynamic target, $k$ 는 그 target 으로 접근하는 kinetic mobility 다.
이 구분이 target 과 mobility 를 동시에 자유 fitting 하는 double counting 을 막는다.
\begin{boundbox}
\textbf{유효범위 (\AL{5}).} 선형 relaxation 은 equilibrium 근처 1차 Taylor 이며 tail 영역
($\theta\approx\theta_\eq$) 에서 일반성을 잃지 않는다 \cite{degrootmazur1962,onsager1931}.
\end{boundbox}

\textbf{Derivation D6 ($q$-coordinate, single-mode kernel).} $v_Q=|I|$, $dq/dt=|I|/Q_\cell$
에서 $d\theta/dq=(Q_\cell k/|I|)\,r$, 그리고
\begin{equation}
\frac{dr}{dq}+\frac{Q_\cell k}{|I|}\,r=\frac{d\theta_\eq}{dq}.
\label{eq:r_ode}
\end{equation}
post-peak 에서 $d\theta_\eq/dq\simeq0$ 이면
\begin{equation}
\boxed{\;r(q)\simeq r(q_a)\exp\!\Big[-\frac{q-q_a}{L_q}\Big],\qquad L_q=\frac{|I|}{Q_\cell\,k}=\frac{v_q}{k}.\;}
\label{eq:single_kernel}
\end{equation}
이것은 관측 tail 전체가 단일 exponential 이라는 뜻이 \emph{아니다} — \textbf{하나의 local
mode 가 만드는 exponential kernel} 이다. 차원: $L_q$ 은 무차원 ($q$ 단위). $|I|/(Q_\cell k)=$
A/(C$\cdot$1/s)$=$무차원. dimensional charge axis 로 쓰면 $L_Q=Q_\cell L_q=v_Q/k$ 이다.

% ====================================================================
\section{★ Barrier distribution $\to$ relaxation-length spectrum}\label{sec:spectrum}
% ====================================================================
\textbf{물리.} 실제 electrode 에는 단일 $G$ 가 아니라, particle size·local staging
environment·domain boundary·defect·stress·surface·transport 차이로 인한 \emph{barrier
distribution} $\rho_G(G;T,\psi)$ 가 있다. 그러나 관측 tail 의 shape 은 $\rho_G$ 와 같지
않다 — $G$ 가 먼저 rate 로, 다시 length 로 변환되기 때문이다.

\textbf{Derivation D7 (barrier$\to$length 지수 매핑).} 식~\eqref{eq:rate} 와 $L_q=|I|/(Q_\cell k)$
에서
\begin{equation}
\boxed{\;L_q(G,T,\psi)=\frac{|I|}{Q_\cell\,k_0(T)}\exp\!\Big[\frac{G-W_\psi}{RT}\Big].\;}
\label{eq:L_of_G}
\end{equation}
$G$ 의 작은 변화가 $L_q$ 에서 \emph{지수적으로 확대}된다 — 이것이 stretched tail 의 핵심
근원이다 (long tail 은 $\rho_G$ 자체가 길어서가 아니라 barrier$\to$length 매핑이 지수이기
때문일 수 있다). 역변환 $G(L_q)=W_\psi+RT\ln[Q_\cell k_0 L_q/|I|]$, Jacobian $|dG/dL_q|=RT/L_q$.
따라서 relaxation-length spectrum:
\begin{equation}
A_L(\lambda;T,\psi)=\rho_G\big(G(\lambda,T,\psi);T,\psi\big)\,\frac{RT}{\lambda}\,A_0(\lambda;T,\psi),
\label{eq:spectrum}
\end{equation}
$A_0$ 는 각 mode 의 initial residual amplitude·population weight·accessibility (무차원
per-mode factor). $\lambda=L_q$ 는 무차원 integration variable 이며, $A_L(\lambda)$ 는
$\lambda$ 에 대한 \emph{density} 다 ($\int A_L(\lambda)\,d\lambda=$ 총 weight).
\begin{boundbox}
\textbf{Non-uniqueness (\AL{6}).} barrier distribution $\to$ stretched kinetics 의 역매핑은
\emph{유일하지 않다} (Plonka \cite{plonka2001}; retrapping 등 다른 메커니즘도 동일 stretched
tail 재현). 따라서 본 Chapter 는 $\rho_G$ 를 관측 tail 에서 \emph{역산하지 않고},
주어진 (또는 독립 측정한) $\rho_G$ 로부터 tail 을 \emph{forward 예측}하는 데에만 쓴다
(\S\ref{sec:falsify}).
\end{boundbox}

% ====================================================================
\section{★ Kernel integral for the observed tail}\label{sec:kernel}
% ====================================================================
전체 phase progress 의 post-peak tail 기여는 local exponential kernel 의 mixture:
\begin{equation}
\boxed{\;\frac{d\Theta_\mathrm{tail}}{dq}\simeq\int_0^\infty A_L(\lambda;T,\psi)\,\frac1\lambda\,
\exp\!\Big[-\frac{q-q_a}{\lambda}\Big]\,d\lambda,\qquad \lambda=L_q.\;}
\label{eq:kernel_integral}
\end{equation}
이것이 본 Chapter 의 중심식이다. single-
mode 식~\eqref{eq:single_kernel} 은 kernel 하나일 뿐이며, 관측 stretched tail 은 spectrum
$A_L$ 의 large-$L_q$ 영역이 결정한다. 식~\eqref{eq:kernel_integral} 은 Laplace-type mixture
이므로 $\rho_G$ 의 형태에 따라 단일 지수~stretched~power-law 사이를 연속적으로 잇는다
(smooth). 일반적인 두 시점 kernel 로 쓰면
\begin{equation}
\mathcal K(q,q')=\int_0^\infty A_L(\lambda;T,\psi)\,\frac1\lambda\,\exp\!\Big[-\frac{q-q'}{\lambda}\Big]\,d\lambda,
\label{eq:K_def}
\end{equation}
이고 식~\eqref{eq:kernel_integral} 은 $\mathcal K(q,q_a)$ 에 해당한다 (이 $\mathcal K$ 가
\S\ref{sec:closure} 의 closure 에 들어가는 kernel).

\textbf{상태 주의 (\AL{6}).} 본 spectrum/kernel 관점은 Assumption Ledger 에서 BOUNDED
이며, 식~\eqref{eq:kernel_integral} 을 본 Chapter 의 중심으로 쓰는 것은 \emph{spectrum
hypothesis 가 \S\ref{sec:falsify} 의 N3/N4 falsification criteria 를 satisfy 한다는 전제} 하에서다
(single-mode 식~\eqref{eq:single_kernel} 은 $A_L(\lambda)=\delta(\lambda-L_{q,0})$ 특수해로 항상 포함).

% ====================================================================
\section{Self-consistent integral 의 closure hierarchy — Plan A / Plan B}\label{sec:closure}
% ====================================================================
\textbf{왜 필요한가.} 식~\eqref{eq:kernel_integral} 은 $A_L$ 을 알면 forward 평가 가능하지만,
$\theta_j$ 와 $V_n$ 이 charge balance (식~\eqref{eq:charge_balance}) 로 결합되어 있어 일반적으로
$\Theta$ 가 적분 안팎에 동시에 나타나는 \emph{self-consistent integral equation} 이 된다.
시간 또는 charge coordinate 적분형은 schematic하게
\begin{equation}
\Theta(q)=\Theta_0+\int_0^q \mathcal K(q,q')\,[\Theta_\eq(q')-\Theta(q')]\,dq'
\label{eq:volterra_self}
\end{equation}
처럼 쓸 수 있다. 여기서 $\mathcal K(q,q')$ 는 식~\eqref{eq:K_def} 의 spectrum kernel 이다.
적분 상한이 $q$ 이므로 원형은 causal memory 를 갖는 Volterra-type second-kind structure 이다.
또한 $\mathcal K$, $\Theta_\eq$, $k$, $A_L$ 은 $V_n$ 에 의존하고, $V_n$ 은 charge balance 로
$\Theta$ 와 다시 연결된다. 이 self-consistency 는 논리 결함이 아니라 graphite ICA 문제의
자연스러운 implicit closure 이다.

\subsection{Closure layer 와 physical core 의 분리}\label{sec:closure_core}
본 장에서 반드시 보존해야 하는 core 는 mathematical shortcut 이 아니라 다음 물리 chain 이다.
\[
\Delta G_\eff \rightarrow k \rightarrow L_q \rightarrow A_L(\lambda;T,\psi) \rightarrow
\int_0^\infty A_L(\lambda;T,\psi)\lambda^{-1}
\exp[-(q-q_a)/\lambda]\,d\lambda \rightarrow \text{ICA tail}.
\]
따라서 closure layer 는 두 가지 역할로 분리한다.
\begin{enumerate}[label=\textbf{I\arabic*.}]
\item \textbf{Physical tail formation.} Activation barrier distribution 이 rate constant
  distribution 으로, 다시 relaxation-length spectrum 으로 바뀌고, 그 spectrum 이 local
  exponential kernel 을 중첩해 observed tail 을 만든다.
\item \textbf{Analytic closure.} $V_n$, $\Theta_\eq$, $k$, $A_L$, $\Theta$ 가 서로 의존하므로,
  paper 또는 later fitting expression 에서 self-consistent integral 을 얼마나 간단한 analytic
  expression 으로 닫을 수 있는지 검토한다.
\end{enumerate}
I1이 Chapter 1의 load-bearing physical logic 이다. I2는 analytic convenience 이며, 특정 closure
method가 성립하지 않아도 I1은 무너지지 않아야 한다. 이 원칙 때문에 본 장은 Plan A와 Plan B를
병기한다.

\begin{center}
\fbox{\parbox{0.92\textwidth}{
\textbf{Plan A.} Fredholm / propagator / ratio-substitution closure 가 현재 graphite ICA
integral structure 에 정합함을 보이면, Refs~\cite{lee2011,son2013} method 를 candidate analytic
closure 로 사용한다.

\textbf{Plan B.} 그 정합성이 증명되지 않거나 Chapter 1 본문에서 analytic closure 를 쓰기 어렵다면,
Refs 6/7 closure 를 핵심식으로 채택하지 않고 local residual ODE + relaxation-length spectrum
kernel integral 을 conservative theoretical formulation 으로 유지한다.
}}
\end{center}

\subsection{Plan A: Refs 6/7 as candidate analytic closure}\label{sec:track1}
사용자의 JCP 2017 paper에서 refs. 6 and 7은 Fredholm integral equation of the second kind를
효율적으로 다루는 propagator / ratio-substitution 계열의 source로 쓰인 것으로 확인되어 있다
\cite{lee2011,son2013}. 이 method가 주는 중요한 guardrail은 다음이다.
\[
\text{self-referential integral equation}
\rightarrow
\text{unknown-function ratio 를 명시}
\rightarrow
\text{controlled approximation of that ratio}
\rightarrow
\text{limiting-case or independent check}.
\]
즉 unknown function 을 조용히 cancel하거나 circular substitution 으로 없애면 안 된다. 어떤 ratio를
어떤 reference solution으로 닫는지, 그 approximation이 어떤 limit에서 유효한지 밝혀야 한다.

Plan A가 Chapter 1의 analytic closure로 승격되려면 최소한 다음 criteria 를 satisfy 해야 한다.
\begin{enumerate}[label=\textbf{M\arabic*.}]
\item Charge-balance feedback 과 barrier spectrum 을 fixed-domain Fredholm-type second-kind
  equation 으로 재정식화할 수 있어야 한다.
\item Unknown function 이 integral 안팎에 같은 function class 로 나타나야 한다.
\item Ratio-substitution ansatz 가 $\theta_j(s)/\theta_j(q)$, $r_j(s)/r_j(q)$, 또는
  $\Theta(s)/\Theta(q)$ 중 어느 대상에 적용되는지 명시되어야 한다.
\item Single-mode limit 에서 식~\eqref{eq:single_kernel} 의 exponential kernel 로 돌아와야 한다.
\item Closure expression 이 Plan B의 local ODE limit 과 같은 limiting behavior 를 가져야 한다.
\end{enumerate}

이 조건이 성립할 때만, 어떤 effective unknown $Y(q)$에 대해
\begin{equation}
Y(q)=Y_0+\int_{\mathcal D}K(q,s)\,\Phi[Y](s)\,\dd s
\label{eq:closure_general}
\end{equation}
를
\begin{equation}
Y_{\mathrm A}(q)=
\mathcal C_{\mathrm{ratio}}\!\left[K,\Phi,Y^{\mathrm{ref}}\right](q)
\label{eq:plan_a_closure}
\end{equation}
와 같은 candidate analytic closure 로 쓸 수 있다. 여기서 $\mathcal C_{\mathrm{ratio}}$ 는
specific ratio-closure operator 이고, 그 detailed form은 graphite ICA의 integral equation class와
variable mapping이 확정된 뒤에만 본문식으로 승격한다. 따라서 Refs 6/7 은 battery-specific
physical assumption 이 아니라 self-consistent integral을 닫는 mathematical method candidate다.

\subsection{Plan B: conservative theoretical formulation}\label{sec:track2}
Plan A의 Fredholm 정식화가 성립하지 않아도 Chapter 1의 physical core는 유지된다. Plan B는
분석상 예비 경로나 계산 절차가 아니라, Chapter 1이 항상 보존해야 하는 논리적 기준식이다.
그 구성은 다음 세 식이다.
\begin{align}
\frac{\dd r_j}{\dd q}+\frac{Q_\cell k_j}{|I|}r_j
&=
\frac{\dd\theta_{\eq,j}}{\dd q},
\label{eq:planb_residual}
\\
L_q(G,T,\psi)
&=
\frac{|I|}{Q_\cell k_0(T)}
\exp\!\left[\frac{G-W_\psi}{RT}\right],
\label{eq:planb_length}
\\
\mathcal T_j(q;T,\psi)
&=
\int_0^\infty
A_{L,j}(\lambda;T,\psi)\frac{1}{\lambda}
\exp\!\left[-\frac{q-q_{a,j}}{\lambda}\right]\dd \lambda,
\qquad \lambda=L_q.
\label{eq:planb_tail}
\end{align}
첫 번째 식은 $q$-coordinate 에서 local state 가 equilibrium target 을 finite rate 로 따라가는 residual equation 이다.
두 번째 식은 activation barrier 와 potential assistance 가 normalized relaxation length 로 바뀌는 exponential
mapping 이다. 세 번째 식은 여러 local kernel 의 spectrum mixture 이다. 이 세 식만으로도
\emph{low $T$에서 large-$L_q$ weight가 커져 tail이 길어지고, positive potential assistance가
short-$L_q$ shift를 만들어 tail을 줄인다}는 Chapter 1의 핵심 설명은 완결된다.

Plan B가 compact analytic expression 을 주지 않는다고 해서 물리 논리가 부족한 것은 아니다. Chapter 1에서
핵심은 state/target/rate/spectrum 의 logical separation 이다. Plan B는 그
separation을 보존하는 baseline theory 이며, Plan A는 이 baseline 위에 얹히는 candidate analytic
compression 이다.

\subsection{Canonical decision rule}\label{sec:closure_decision}
본 장의 canonical rule은 다음과 같다.
\begin{enumerate}[label=\textbf{R\arabic*.}]
\item Plan A는 validated analytic closure가 될 수 있지만, 검증 전에는 core assumption이 아니다.
\item Plan B는 항상 남는 conservative theoretical base다.
\item 논문 본문에서 Plan A를 사용하더라도 appendix 또는 methods section에는 Plan B derivation과
  limiting-case consistency를 함께 제시한다.
\item Plan A와 Plan B가 같은 limiting cases를 주지 않으면, Chapter 1의 core 논리로는 Plan B
  formulation만 채택한다.
\item 이 dual hierarchy는 계산 절차가 아니라 논리적 책임 분리다.
\end{enumerate}
\section{ICA mapping}\label{sec:ica}
% ====================================================================
background storage $dQ_\bg=C_\bg(V_n,T)\,dV_n$ 와 phase progress $Q_p\,d\Theta$ 의 합은
\[
dQ_\ext=C_\bg(V_n,T)\,dV_n+Q_p\,d\Theta,
\qquad dQ_\ext=Q_\cell\,dq
\]
이다. 식~\eqref{eq:charge_balance_differential} 과 동치로 정리하면
\begin{equation}
\boxed{\;\frac{dQ_\ext}{dV_n}=Q_\cell\frac{dq}{dV_n}
=
\frac{C_\bg(V_n,T)}
{1-(Q_p/Q_\cell)\,(d\Theta/dq)}.\;}
\label{eq:ica}
\end{equation}
식~\eqref{eq:kernel_integral} 의 tail 기여가 $d\Theta/dq$ 에 들어가 ICA tail 로 나타난다.
dimensional charge coordinate 로 바꾸어 읽어야 할 때는 $d\Theta/dQ_\ext=(1/Q_\cell)d\Theta/dq$ 와
$L_Q=Q_\cell L_q$ 를 사용한다.

\textbf{측정 축 bridge (두 단계 구분).} 이론의 primary potential 은 internal anode
$V_n$ (charge-balance root, 식~\eqref{eq:charge_balance}) 이고 rate/barrier 인수는
$V_{n,\dimedrive}=V_n$ (기본). 측정 축으로 가는 데 \emph{서로 다른 두 양}을 구분해야 한다.
\begin{enumerate}[label=(\roman*)]
\item \textbf{Anode apparent potential} $V_{n,\app}$ — anode 자신의 polarization 을 더한
  양 ($\sim0.1$ V scale): $V_{n,\app}=V_n+s_I|I|R_n$, $R_n$ 은 anode lumped polarization
  계수 (charge-transfer+transport+film 합, \S\ref{sec:notation}).
\item \textbf{Full-cell terminal voltage} $V_{\cell}$ — anode·cathode·ohmic 의 차
  ($\sim3$--$4$ V scale), $V_{n,\app}$ 와 \emph{다른} 물리량:
  \begin{equation}
  V_{\cell}(q)=V_p(q)-V_{n,\app}(q)-I R_{\Omega,\mathrm{sep}}(q,T),
  \label{eq:vapp_bridge}
  \end{equation}
  ($V_p$ cathode apparent, $R_{\Omega,\mathrm{sep}}$ separator/electrolyte ohmic; anode
  의 $\eta_{\mathrm{ct}},\eta_{\mathrm{tr}}$ 는 이미 $V_{n,\app}$ 의 $R_n$ 에 흡수됨 —
  중복 계상 금지).
\end{enumerate}
half-cell 측정이면 $V_{n,\app}$ 를 직접 보고, full-cell 이면 식~\eqref{eq:vapp_bridge} 로
anode 기여를 분리해야 한다. voltage-axis tail scale 은 normalized coordinate 에서
$L_\varphi\simeq|dV_n/dq|_{q_a}L_q$ 이며, dimensional coordinate 로는
$L_\varphi\simeq|dV_n/dQ_\ext|_{q_a}L_Q$ 이다.
이 구분(특히 $V_{n,\dimedrive}$ 의 potential assistance vs $R_n$ 의 polarization shift)을
놓치면 \S\ref{sec:falsify} 의 double counting 위험이 있다.

% ====================================================================
\section{★ Spectrum 의 $T/\psi$ 의존 = stretched tail (O2)}\label{sec:tailT}
% ====================================================================
\textbf{Derivation D9 (novel; \AL{8}).} 식~\eqref{eq:L_of_G} 에서:
\begin{itemize}
\item \textbf{Low $T$}: $\exp[(G-W_\psi)/RT]$ 가 커지고 $k_0(T)$ 가 작아져 같은 $G-W_\psi$
  에 대해 $L_q$ 이 커진다 $\Rightarrow$ $A_L$ weight 가 large-$L_q$ 로 이동 $\Rightarrow$
  식~\eqref{eq:kernel_integral} 의 long-$q$ tail 이 길어진다 (저온 긴 tail).
\item \textbf{High $T$}: $RT$ 증가로 $L_q$ 이 작아져 $A_L$ 이 short-$L_q$ 로 이동 $\Rightarrow$
  tail 이 짧게 끝난다.
\item \textbf{Potential assistance}: $W_\psi=\chi_j\mathcal A_j$ 증가 시 $L_q$ 감소,
  \begin{equation}
  \frac{\partial\ln L_q}{\partial V_{n,\dimedrive}}=-\frac{\chi_j s_{\phi,j}F}{RT}<0\ (\text{forward}),
  \label{eq:psi_shift}
  \end{equation}
  즉 potential 이 spectrum 을 short-$L_q$ 로 shift $\Rightarrow$ tail 단축. 이는 state jump
  가 아니라 \emph{spectrum shift} 다.
\end{itemize}
\textbf{Equilibrium-width 와의 구별.} homogeneous lattice-gas width 는 $w_j=RT/F$ 라 저온서
오히려 좁아진다. 따라서 ``저온 긴 tail''은 homogeneous-equilibrium 효과로 설명되지 않으며
\emph{kinetic spectrum shift} 가 필요하다 (단, heterogeneous disorder width 는 $T$-약의존
이므로 별도 — \S\ref{sec:falsify} lineshape test). 본 절은 kinetic-spectrum 메커니즘이
(O2) 와 \emph{강하게 시사(consistent)} 됨을 보이며, \emph{확정}은 falsification criteria 만족 전제.

% ====================================================================
\section{★ Falsification 과 non-uniqueness 처리}\label{sec:falsify}
% ====================================================================
\textbf{왜 필요한가.} 관측 tail 을 barrier-spectrum 에 귀속하려면 competing source 를
배제해야 한다 (\AL{10}; ICA 가 rate 로 broadening 된다는 것은 알려져 있으므로
\cite{dubarry2022,fly2020} 단순 rate-broadening 과 구별이 핵심).
\begin{longtable}{p{0.24\textwidth} p{0.40\textwidth} p{0.30\textwidth}}
\toprule Competing source & tail 기여 & 판별 \\ \midrule \endhead
Equilibrium width & homogeneous/disorder broadening & 저속 OCV peak near-peak lineshape; 저온서 좁아지면 kinetic 아님 (homogeneous 한정) \\
Polarization $R_n$ & $V_{n,\app}$ shift & pulse/current-interrupt/다중 $|I|$ \\
Transport/diffusion & 저온 느린 확산 lag & GITT relaxation signature; $D(T)$ Arrhenius vs barrier slope \\
Barrier distribution $\rho_G$ \cite{plonka2001} & spectrum kernel & \emph{비유일} 역매핑 — forward prediction 만; 역산 금지 \\
\bottomrule
\end{longtable}

\textbf{★ 비퇴화 discriminator.} tail length scaling 만으로는 부족하다: $L_q\propto|I|$ 는
transport lag 도 공유한다 (퇴화). barrier-spectrum 의 고유 signature 는 식~\eqref{eq:L_of_G},
\eqref{eq:psi_shift} 의 \emph{potential-dependent} 항이다. 즉 (a) tail length 의 Arrhenius
slope 자체가 $(V_{n,\dimedrive}-U_j)$ 에 의존하고, (b) potential/$|I|$ 증가 시 spectrum 이
$\chi_j$ 를 통해 short-$L_q$ 로 shift 한다. transport 와 \emph{단순} polarization 은 이런
potential-dependent slope 를 갖지 않는다.

\textbf{★ 가장 약한 지점 — $\chi_j$ vs charge-transfer $\eta_{\mathrm{ct}}$ 의 co-linearity.}
주의: charge-transfer polarization $\eta_{\mathrm{ct}}$ 도 Butler--Volmer 라 potential 에
지수적이고 $R_n$ 은 $q$(따라서 potential) 의존이므로, $\chi_j$ barrier-lowering slope 와
$\eta_{\mathrm{ct}}$ Tafel slope 는 \emph{단일 electrode 데이터에서 co-linear (퇴화)} 일 수
있다. 따라서 (b) 신호만으로 $\chi_j$ 를 확정할 수 없다. 이 둘을 가르려면 \S\ref{sec:fitting}
의 다중 $T\times|I|$ + GITT/pulse 설계가 필수다 (GITT 의 instantaneous step 이 $\eta_{\mathrm{ct}}$
를, slow relaxation 이 barrier-spectrum 을 분리). 즉 $\chi_j$ 귀속은 $\eta_{\mathrm{ct}}$ 를
독립 분리한 \emph{후}에만 성립한다.

\textbf{Null-result 규칙.}
\begin{enumerate}[label=(N\arabic*)]
\item tail length 의 $\ln$ vs $1/T$ slope 가 독립 측정 $E_D$ 와 일치하고 potential 의존이
  없으면 $\Rightarrow$ barrier-spectrum 귀속 기각 (transport-dominant).
\item potential/$|I|$ 증가 시 spectrum/peak shift 가 없으면 $\Rightarrow$ $\chi_j$ 항 부재
  $\Rightarrow$ potential-lowering 가설 기각.
\item 저속 OCV near-peak lineshape 이 저온서 좁아지지 않으면 $\Rightarrow$ heterogeneous
  disorder 기여 존재 — kinetic 단독 귀속 기각, disorder 와 분리.
\item $\rho_G$ 를 관측 tail 에서 역산하지 않는다 (non-unique); $\rho_G$ 는 독립 근거
  (입자분포·EIS·계산) 로 주고 tail 을 forward 예측해 비교만 한다.
\end{enumerate}

% ====================================================================
\section{Later fitting interface, identifiability, validation}\label{sec:fitting}
% ====================================================================
\textbf{Deliverable.} Chapter 1이 later quantitative use 로 넘기는 기준식은 먼저 Plan B
baseline 이다:
\[
k(G,T,\psi)\ \xrightarrow{\eqref{eq:rate}}\ 
L_q(G,T,\psi)\ \xrightarrow{\eqref{eq:L_of_G}}\ 
A_L(\lambda;T,\psi)\ \xrightarrow{\eqref{eq:spectrum}}\ 
\mathcal T_j(q;T,\psi)\ \xrightarrow{\eqref{eq:planb_tail}}\ 
\frac{d\Theta_\mathrm{tail}}{dq}\ \xrightarrow{\eqref{eq:ica}}\ 
\frac{dQ_\ext}{dV_n}.
\]
Plan A는 Refs 6/7 ratio-substitution structure 가 M1--M5 와 limiting-case check 를 satisfy 할 때,
self-consistent layer 를 더 compact 하게 쓰는 candidate analytic compression 이다. 따라서 later
fitting expression 은 Plan B baseline 을 항상 보존하고, Plan A는 검증 후 선택적으로 치환한다.
\textbf{유효범위}: $\Delta G_\eff\ge0$ (Marcus bound), tail 영역
$V_{n,\dimedrive}\approx U_j+O(w_j)$.

\textbf{Parameter / 실험 매핑 (\AL{10}).}
\begin{longtable}{p{0.26\textwidth} p{0.30\textwidth} p{0.38\textwidth}}
\toprule parameter & 의미 & 결정 실험 \\ \midrule \endhead
$U_j(T),w_j(T),Q_p$ & equilibrium target & 저속 OCV / GITT 준평형 \\
$\Delta H_a,\Delta S_a$ & activation & 다중 $T$ Arrhenius \\
$\chi_j$ (혹은 $\Lambda_\psi$) & potential lowering & 다중 $|I|$/potential shift \\
$\rho_G$ (외부 입력) & barrier spectrum & 독립 근거(입자분포/EIS/계산)로 \emph{주어짐} — tail 역산 금지 (forward only) \\
$k_0(T)$ & prefactor & Arrhenius intercept \\
\bottomrule
\end{longtable}
\textbf{Non-identifiability 경고.} $R_n,\rho_G,w_j,k_0$ 는 모두 broadening 에 기여하므로
단일 $T$/제한된 $|I|$ 로 동시 분리 불가. 독립 실험(pulse/GITT/EIS, 다중 $T\times|I|$) 으로
제약 필수. quantitative fitting implementation 은 Chapter 1 범위 외 (P1).

% ====================================================================
\section{종합·validation status·참고문헌}\label{sec:summary}
% ====================================================================
\subsection{결과 요약}
\begin{enumerate}[label=\textbf{(R\arabic*)}]
\item charge conservation 으로 $V_n$ implicit (S1, \AL{9}).
\item equilibrium target 일반형(smooth); lattice-gas grounded 특수해, 형태 data 결정 (S2, \AL{3a}).
\item ★ effective barrier $\Delta G_\eff=\Delta G_a-\chi\mathcal A=G-W_\psi$, Marcus-bounded (S5, \AL{2}).
\item single-mode relaxation + exponential kernel $L_q=|I|/(Q_\cell k)$ and dimensional conversion $L_Q=Q_\cell L_q$ (S6/S7, \AL{1},\AL{5}).
\item ★ barrier distribution $\to$ 지수 매핑 $\to$ relaxation-length spectrum $A_L$ (S8, \AL{6}).
\item ★ 관측 tail = spectrum kernel integral (S9); closure hierarchy = \textbf{Plan B baseline (local residual ODE + spectrum kernel, 항상 보존) + Plan A candidate analytic compression (Refs 6/7, M1-M5+$\varepsilon$ gated) +
  R1-R5 governance} (S10, \AL{7}, VG-2).
\item ★ stretched tail 의 $T/\psi$ 의존 = spectrum shift (novel; S12, \AL{8}); homogeneous-equilibrium 으로 설명 불가, kinetic-spectrum 강하게 시사 (확정은 falsification criteria 만족 전제).
\item ★ falsification protocol + 비퇴화 discriminator + $\rho_G$ non-unique inversion 회피 (S13, \AL{10}).
\item later fitting interface + identifiability + validation (S14).
\end{enumerate}

\subsection{Assumption and validation status}
\begin{longtable}{p{0.58\textwidth} p{0.34\textwidth}}
\toprule 검증 항목 & 현재 상태 \\ \midrule \endhead
모든 본문 assumption 이 AL-\# 에 연결됨 & documented \\
FLAGGED(AL-3b erf, AL-7 applicability, AL-8 tail) 를 established result 로 쓰지 않음 & controlled as hypothesis \\
BOUNDED(AL-2 Marcus, AL-5 near-equilibrium, AL-6 non-unique inversion) 의 유효범위 병기 & bounded with stated scope \\
step-function/$0\!\to\!1$ 급점프 정의식 부재 & excluded by smooth-only convention \\
activation barrier 유지 & retained as core mechanism (\S\ref{sec:barrier}) \\
``가우시안''을 equilibrium shape 확정으로 쓰지 않음 & treated as descriptive peak-shape example (\S\ref{sec:equilibrium}) \\
tail 귀속을 established result 가 아니라 consistency + falsification protocol 로 둠 & stated as hypothesis with tests (\S\ref{sec:tailT},\ref{sec:falsify}) \\
$\rho_G$ non-unique inversion 회피 & forward model only (\S\ref{sec:spectrum},\ref{sec:falsify}) \\
사용자 PhD Refs 6/7 = Plan A analytic closure candidate & gated by M1-M5+$\varepsilon$; Plan B core preserved (\S\ref{sec:closure}) \\
\bottomrule
\end{longtable}

\subsection{Chapter 2+ 로 전달}
single-mode relaxation, spectrum kernel, $\Delta G_\eff$ 가 reversible/irreversible heat,
$dV/dT$, Butler--Volmer kinetics, 충방전 hysteresis(branch-dependent spectrum), full-cell
의 input (별도 roadmap).

\begin{thebibliography}{99}
\bibitem{eyring1935} H. Eyring, ``The activated complex in chemical reactions,''
  \emph{J. Chem. Phys.} \textbf{3}, 107 (1935). DOI: 10.1063/1.1749604.
\bibitem{eyringchemrev1935} S. Glasstone, K. J. Laidler, H. Eyring, ``The activated
  complex and the absolute rate of chemical reactions,'' \emph{Chem. Rev.} \textbf{17},
  301 (1935). DOI: 10.1021/cr60056a006.
\bibitem{evanspolanyi1938} M. G. Evans, M. Polanyi, ``Inertia and driving force of
  chemical reactions,'' \emph{Trans. Faraday Soc.} \textbf{34}, 11 (1938). DOI:
  10.1039/TF9383400011.
\bibitem{marcus1956} R. A. Marcus, ``On the theory of oxidation--reduction reactions
  involving electron transfer,'' \emph{J. Chem. Phys.} \textbf{24}, 966 (1956). DOI:
  10.1063/1.1742723.
\bibitem{bardfaulkner} A. J. Bard, L. R. Faulkner, \emph{Electrochemical Methods},
  Wiley (Ch.~3).
\bibitem{mckinnon1983} W. R. McKinnon, R. R. Haering, ``Physical mechanisms of
  intercalation,'' \emph{Modern Aspects of Electrochemistry} \textbf{15}, 235 (1983).
\bibitem{bazant2013} M. Z. Bazant, ``Theory of chemical kinetics and charge transfer
  based on nonequilibrium thermodynamics,'' \emph{Acc. Chem. Res.} \textbf{46}, 1144
  (2013). DOI: 10.1021/ar300145c.
\bibitem{dahn1991} J. R. Dahn, ``Phase diagram of Li$_x$C$_6$,'' \emph{Phys. Rev. B}
  \textbf{44}, 9170 (1991). DOI: 10.1103/PhysRevB.44.9170.
\bibitem{ohzuku1993} T. Ohzuku, Y. Iwakoshi, K. Sawai, \emph{J. Electrochem. Soc.}
  \textbf{140}, 2490 (1993). DOI: 10.1149/1.2220849.
\bibitem{funabiki1999} A. Funabiki, M. Inaba, T. Abe, Z. Ogumi, ``Nucleation and
  phase-boundary movement upon stage transformation in lithium--graphite intercalation
  compounds,'' \emph{Electrochim. Acta} \textbf{45}, 865 (1999). DOI:
  10.1016/S0013-4686(99)00290-X.
\bibitem{degrootmazur1962} S. R. de Groot, P. Mazur, \emph{Non-Equilibrium
  Thermodynamics}, North-Holland (1962); Dover (1984), Ch.~X.
\bibitem{onsager1931} L. Onsager, ``Reciprocal relations in irreversible processes,''
  \emph{Phys. Rev.} \textbf{37}, 405 (1931).
\bibitem{plonka2001} A. Plonka, \emph{Dispersive Kinetics}, Springer (2001).
\bibitem{lee2011} S. Lee, C. Y. Son, J. Sung, S. H. Chong, ``Communication: Propagator
  for diffusive dynamics of an interacting molecular pair,'' \emph{J. Chem. Phys.}
  \textbf{134}, 121102 (2011). DOI: 10.1063/1.3565476.
\bibitem{son2013} C. Y. Son, J. Kim, J.-H. Kim, J. S. Kim, S. Lee, ``An accurate
  expression for the rates of diffusion-influenced bimolecular reactions with long-range
  reactivity,'' \emph{J. Chem. Phys.} \textbf{138}, 164123 (2013). DOI: 10.1063/1.4802584.
\bibitem{doyle1993} M. Doyle, T. F. Fuller, J. Newman, ``Modeling of galvanostatic
  charge and discharge,'' \emph{J. Electrochem. Soc.} \textbf{140}, 1526 (1993). DOI:
  10.1149/1.2221597.
\bibitem{newman} J. Newman, K. E. Thomas-Alyea, \emph{Electrochemical Systems}, 3rd ed.,
  Wiley (2004).
\bibitem{baessler1993} H. Bässler, ``Charge transport in disordered organic
  photoconductors,'' \emph{Phys. Status Solidi B} \textbf{175}, 15 (1993). DOI:
  10.1002/pssb.2221750102.
\bibitem{dubarry2022} M. Dubarry, D. Anseán, ``Best practices for incremental capacity
  analysis,'' \emph{Front. Energy Res.} \textbf{10}, 1023555 (2022). DOI:
  10.3389/fenrg.2022.1023555.
\bibitem{fly2020} A. Fly, R. Chen, ``Rate dependency of incremental capacity analysis
  (dQ/dV),'' \emph{J. Energy Storage} \textbf{29}, 101329 (2020). DOI:
  10.1016/j.est.2020.101329.
\end{thebibliography}

\end{document}


